\documentclass[11pt]{article}
\newcommand{\ListaTitulo}{Gabarito --- Lista 09}
% Preâmbulo compartilhado — MATD48, Listas de Exercícios 2026
% Usado por Lista{NN}.tex e Gabarito{NN}.tex via % Preâmbulo compartilhado — MATD48, Listas de Exercícios 2026
% Usado por Lista{NN}.tex e Gabarito{NN}.tex via % Preâmbulo compartilhado — MATD48, Listas de Exercícios 2026
% Usado por Lista{NN}.tex e Gabarito{NN}.tex via \input{preamble}

\usepackage[utf8]{inputenc}
\usepackage[T1]{fontenc}
\usepackage[brazil]{babel}
\usepackage{fullpage}
\usepackage{amsmath,amssymb}
\usepackage{enumitem}
\usepackage{xcolor}
\usepackage{fancyhdr}
\usepackage{titlesec}
\usepackage{listings}
\usepackage{hyperref}
\usepackage{booktabs}
\usepackage{graphicx}

\definecolor{matd48azul}{HTML}{0B4F6C}
\definecolor{matd48verde}{HTML}{2E8B57}

\titleformat{\section}{\normalfont\Large\bfseries\color{matd48azul}}{\thesection}{1em}{}
\titleformat{\subsection}{\normalfont\large\bfseries\color{matd48azul}}{\thesubsection}{1em}{}

\providecommand{\ListaTitulo}{Lista de Exercícios}

\setlength{\headheight}{14pt}
\pagestyle{fancy}
\fancyhf{}
\lhead{\small MATD48 --- Planejamento de Experimentos A}
\rhead{\small \ListaTitulo}
\cfoot{\thepage}
\renewcommand{\headrulewidth}{0.4pt}

\lstset{
  basicstyle=\ttfamily\small,
  breaklines=true,
  frame=single,
  columns=fullflexible,
  backgroundcolor=\color{gray!8},
  commentstyle=\color{matd48verde},
  keywordstyle=\color{matd48azul}\bfseries,
  language=R,
  extendedchars=true,
  literate=%
    {á}{{\'a}}1 {à}{{\`a}}1 {â}{{\^a}}1 {ã}{{\~a}}1
    {é}{{\'e}}1 {ê}{{\^e}}1 {í}{{\'i}}1 {ó}{{\'o}}1
    {ô}{{\^o}}1 {õ}{{\~o}}1 {ú}{{\'u}}1 {ç}{{\c c}}1
    {Á}{{\'A}}1 {À}{{\`A}}1 {Â}{{\^A}}1 {Ã}{{\~A}}1
    {É}{{\'E}}1 {Ê}{{\^E}}1 {Í}{{\'I}}1 {Ó}{{\'O}}1
    {Ô}{{\^O}}1 {Õ}{{\~O}}1 {Ú}{{\'U}}1 {Ç}{{\c C}}1
}

\hypersetup{colorlinks=true, linkcolor=matd48azul, urlcolor=matd48azul, citecolor=matd48azul}

\newcommand{\cabecalho}[2]{%
  \begin{center}
    {\Large\bfseries\color{matd48azul} #1}\\[0.3em]
    {\large #2}\\[0.3em]
    \rule{\linewidth}{0.6pt}
  \end{center}
  \vspace{0.5em}
}

\newenvironment{questao}[1]{\vspace{1em}\noindent\textbf{Questão #1.}\hspace{0.4em}}{\par}
\newenvironment{solucao}{\vspace{0.3em}\noindent\textit{Solução.}\hspace{0.4em}\color{black!85}}{\par\vspace{0.5em}}


\usepackage[utf8]{inputenc}
\usepackage[T1]{fontenc}
\usepackage[brazil]{babel}
\usepackage{fullpage}
\usepackage{amsmath,amssymb}
\usepackage{enumitem}
\usepackage{xcolor}
\usepackage{fancyhdr}
\usepackage{titlesec}
\usepackage{listings}
\usepackage{hyperref}
\usepackage{booktabs}
\usepackage{graphicx}

\definecolor{matd48azul}{HTML}{0B4F6C}
\definecolor{matd48verde}{HTML}{2E8B57}

\titleformat{\section}{\normalfont\Large\bfseries\color{matd48azul}}{\thesection}{1em}{}
\titleformat{\subsection}{\normalfont\large\bfseries\color{matd48azul}}{\thesubsection}{1em}{}

\providecommand{\ListaTitulo}{Lista de Exercícios}

\setlength{\headheight}{14pt}
\pagestyle{fancy}
\fancyhf{}
\lhead{\small MATD48 --- Planejamento de Experimentos A}
\rhead{\small \ListaTitulo}
\cfoot{\thepage}
\renewcommand{\headrulewidth}{0.4pt}

\lstset{
  basicstyle=\ttfamily\small,
  breaklines=true,
  frame=single,
  columns=fullflexible,
  backgroundcolor=\color{gray!8},
  commentstyle=\color{matd48verde},
  keywordstyle=\color{matd48azul}\bfseries,
  language=R,
  extendedchars=true,
  literate=%
    {á}{{\'a}}1 {à}{{\`a}}1 {â}{{\^a}}1 {ã}{{\~a}}1
    {é}{{\'e}}1 {ê}{{\^e}}1 {í}{{\'i}}1 {ó}{{\'o}}1
    {ô}{{\^o}}1 {õ}{{\~o}}1 {ú}{{\'u}}1 {ç}{{\c c}}1
    {Á}{{\'A}}1 {À}{{\`A}}1 {Â}{{\^A}}1 {Ã}{{\~A}}1
    {É}{{\'E}}1 {Ê}{{\^E}}1 {Í}{{\'I}}1 {Ó}{{\'O}}1
    {Ô}{{\^O}}1 {Õ}{{\~O}}1 {Ú}{{\'U}}1 {Ç}{{\c C}}1
}

\hypersetup{colorlinks=true, linkcolor=matd48azul, urlcolor=matd48azul, citecolor=matd48azul}

\newcommand{\cabecalho}[2]{%
  \begin{center}
    {\Large\bfseries\color{matd48azul} #1}\\[0.3em]
    {\large #2}\\[0.3em]
    \rule{\linewidth}{0.6pt}
  \end{center}
  \vspace{0.5em}
}

\newenvironment{questao}[1]{\vspace{1em}\noindent\textbf{Questão #1.}\hspace{0.4em}}{\par}
\newenvironment{solucao}{\vspace{0.3em}\noindent\textit{Solução.}\hspace{0.4em}\color{black!85}}{\par\vspace{0.5em}}


\usepackage[utf8]{inputenc}
\usepackage[T1]{fontenc}
\usepackage[brazil]{babel}
\usepackage{fullpage}
\usepackage{amsmath,amssymb}
\usepackage{enumitem}
\usepackage{xcolor}
\usepackage{fancyhdr}
\usepackage{titlesec}
\usepackage{listings}
\usepackage{hyperref}
\usepackage{booktabs}
\usepackage{graphicx}

\definecolor{matd48azul}{HTML}{0B4F6C}
\definecolor{matd48verde}{HTML}{2E8B57}

\titleformat{\section}{\normalfont\Large\bfseries\color{matd48azul}}{\thesection}{1em}{}
\titleformat{\subsection}{\normalfont\large\bfseries\color{matd48azul}}{\thesubsection}{1em}{}

\providecommand{\ListaTitulo}{Lista de Exercícios}

\setlength{\headheight}{14pt}
\pagestyle{fancy}
\fancyhf{}
\lhead{\small MATD48 --- Planejamento de Experimentos A}
\rhead{\small \ListaTitulo}
\cfoot{\thepage}
\renewcommand{\headrulewidth}{0.4pt}

\lstset{
  basicstyle=\ttfamily\small,
  breaklines=true,
  frame=single,
  columns=fullflexible,
  backgroundcolor=\color{gray!8},
  commentstyle=\color{matd48verde},
  keywordstyle=\color{matd48azul}\bfseries,
  language=R,
  extendedchars=true,
  literate=%
    {á}{{\'a}}1 {à}{{\`a}}1 {â}{{\^a}}1 {ã}{{\~a}}1
    {é}{{\'e}}1 {ê}{{\^e}}1 {í}{{\'i}}1 {ó}{{\'o}}1
    {ô}{{\^o}}1 {õ}{{\~o}}1 {ú}{{\'u}}1 {ç}{{\c c}}1
    {Á}{{\'A}}1 {À}{{\`A}}1 {Â}{{\^A}}1 {Ã}{{\~A}}1
    {É}{{\'E}}1 {Ê}{{\^E}}1 {Í}{{\'I}}1 {Ó}{{\'O}}1
    {Ô}{{\^O}}1 {Õ}{{\~O}}1 {Ú}{{\'U}}1 {Ç}{{\c C}}1
}

\hypersetup{colorlinks=true, linkcolor=matd48azul, urlcolor=matd48azul, citecolor=matd48azul}

\newcommand{\cabecalho}[2]{%
  \begin{center}
    {\Large\bfseries\color{matd48azul} #1}\\[0.3em]
    {\large #2}\\[0.3em]
    \rule{\linewidth}{0.6pt}
  \end{center}
  \vspace{0.5em}
}

\newenvironment{questao}[1]{\vspace{1em}\noindent\textbf{Questão #1.}\hspace{0.4em}}{\par}
\newenvironment{solucao}{\vspace{0.3em}\noindent\textit{Solução.}\hspace{0.4em}\color{black!85}}{\par\vspace{0.5em}}


\begin{document}

\cabecalho{Gabarito --- Lista de Exercícios 09}{Delineamento em blocos completos aleatorizados (Capítulo 4 / Aula 09)}

\begin{questao}{1} \textbf{(Agricultura --- tabela de ANOVA do DBCA)}
\end{questao}
\begin{solucao}
\begin{enumerate}[label=(\alph*)]
  \item $SQ_{Erro} = SQ_{Total}-SQ_{Trat}-SQ_{Bloco} = 850-320-210 = \textbf{320}$.
    Graus de liberdade: $gl_{Trat}=t-1=3$, $gl_{Bloco}=b-1=4$, $gl_{Erro}=(t-1)(b-1)=12$,
    $gl_{Total}=tb-1=19$. Quadrados médios: $QM_{Trat}=320/3=106{,}67$,
    $QM_{Bloco}=210/4=52{,}5$, $QM_{Erro}=320/12=26{,}67$.
  \item $F_0 = QM_{Trat}/QM_{Erro} = 106{,}67/26{,}67 \approx \textbf{4{,}00}$.
  \item Como $F_0=4{,}00 > F_{0{,}05;3;12}=3{,}49$, \textbf{rejeita-se $H_0$}: há evidência, ao
    nível de $5\%$, de que pelo menos uma variedade difere das demais quanto à produtividade
    média.
\end{enumerate}
\end{solucao}

\begin{questao}{2} \textbf{(Dedução --- decomposição ortogonal da soma de quadrados)}
\end{questao}
\begin{solucao}
\begin{enumerate}[label=(\alph*)]
  \item $\displaystyle\sum_i a_i = \sum_i(\bar Y_{i\cdot}-\bar Y_{\cdot\cdot}) =
    \Big(\sum_i \bar Y_{i\cdot}\Big) - t\bar Y_{\cdot\cdot}$. Como $\bar Y_{\cdot\cdot}$ é, por
    definição, a média das médias de tratamento $\bar Y_{i\cdot}$ (o desenho é balanceado, cada
    tratamento com $b$ observações), $\sum_i \bar Y_{i\cdot} = t\bar Y_{\cdot\cdot}$, logo
    $\sum_i a_i = 0$. O argumento para $\sum_j b_j=0$ é idêntico, trocando tratamento por bloco.
  \item Para $i$ fixo,
    $$
    \sum_j e_{ij} = \sum_j Y_{ij} - b\bar Y_{i\cdot} - \sum_j \bar Y_{\cdot j} + b\bar Y_{\cdot\cdot}.
    $$
    Mas $\sum_j Y_{ij} = b\bar Y_{i\cdot}$ (soma das $b$ observações do tratamento $i$) e
    $\sum_j \bar Y_{\cdot j} = b\bar Y_{\cdot\cdot}$ (pelo item (a), aplicado aos blocos). Logo
    $\sum_j e_{ij} = b\bar Y_{i\cdot} - b\bar Y_{i\cdot} - b\bar Y_{\cdot\cdot} + b\bar
    Y_{\cdot\cdot} = 0$. O argumento para $\sum_i e_{ij}=0$ (com $j$ fixo) é análogo, trocando os
    papéis de $i$ e $j$.
  \item Elevando ao quadrado e somando,
    $$
    \sum_{i,j}(Y_{ij}-\bar Y_{\cdot\cdot})^2 = \sum_{i,j}(a_i+b_j+e_{ij})^2 =
    \sum_{i,j}a_i^2 + \sum_{i,j}b_j^2 + \sum_{i,j}e_{ij}^2
    + 2\sum_{i,j}a_ib_j + 2\sum_{i,j}a_ie_{ij} + 2\sum_{i,j}b_je_{ij}.
    $$
    O primeiro termo cruzado se anula porque $a_i$ não depende de $j$ nem $b_j$ de $i$:
    $\sum_{i,j}a_ib_j = \big(\sum_i a_i\big)\big(\sum_j b_j\big) = 0\times 0 = 0$ pelo item (a).
    O segundo se anula porque, para cada $i$, $a_i$ é constante em $j$ e $\sum_j e_{ij}=0$ (item
    b): $\sum_{i,j}a_ie_{ij} = \sum_i a_i\big(\sum_j e_{ij}\big) = \sum_i a_i\times 0 = 0$.
    Analogamente, $\sum_{i,j}b_je_{ij}=\sum_j b_j\big(\sum_i e_{ij}\big)=0$. Restam apenas os
    três termos quadráticos:
    $$
    \sum_{i,j}(Y_{ij}-\bar Y_{\cdot\cdot})^2 = b\sum_i a_i^2 + t\sum_j b_j^2 + \sum_{i,j}e_{ij}^2
    = SQ_{Trat} + SQ_{Bloco} + SQ_{Erro},
    $$
    em que o fator $b$ multiplicando $\sum_i a_i^2$ (e $t$ multiplicando $\sum_j b_j^2$) surge
    porque $a_i^2$ (resp. $b_j^2$) é somado $b$ vezes (resp. $t$ vezes) — uma vez para cada valor
    do índice do qual não depende. \textbf{Nota:} este é o mesmo fato que, na Seção de notação
    matricial do Capítulo 4, aparece como ortogonalidade entre as colunas (centralizadas) de
    $X_{trat}$ e $X_{bloco}$ na matriz de delineamento $X=[\mathbf1\mid X_{trat}\mid X_{bloco}]$
    — a prova algébrica acima é a mesma ortogonalidade, escrita sem notação matricial.
\end{enumerate}
\end{solucao}

\begin{questao}{3} \textbf{(Psicologia --- blocagem como aleatorização restrita)}
\end{questao}
\begin{solucao}
\begin{enumerate}[label=(\alph*)]
  \item Sejam $Y_u(1)$ e $Y_u(3)$ os resultados potenciais do paciente $u$ sob respiração guiada e
    sob mindfulness breve, respectivamente. O efeito causal médio de interesse é
    $\tau_{13}=\mathbb E\big[Y_u(1)-Y_u(3)\big]$, tomado sobre a população dos pacientes
    recrutados.
  \item A aleatorização — mesmo restrita a cada bloco de gravidade — garante que, \emph{dentro de
    cada bloco}, os pacientes que efetivamente receberam a técnica $i$ são uma amostra
    representativa de como todos os pacientes \emph{daquele bloco} responderiam à técnica $i$: o
    sorteio não depende de nenhuma característica do paciente (observada ou não) que também afete
    a resposta, então não há confusão entre "recebeu a técnica $i$" e "teria respondido melhor de
    qualquer forma". Como isso vale bloco a bloco, e o efeito causal médio geral é uma combinação
    ponderada (pelo tamanho de cada bloco) dos efeitos médios dentro de cada bloco, o estimador
    permanece não viesado para $\tau_{13}$ — a blocagem apenas organiza \emph{onde} a
    aleatorização acontece, sem comprometer a comparabilidade que ela garante.
  \item A blocagem reduz o erro-padrão quando a gravidade basal está de fato correlacionada com a
    resposta (redução de ansiedade) — por exemplo, se pacientes com ansiedade grave tendem a ter
    reduções absolutas maiores (ou menores) do que os leves, independentemente da técnica. Nesse
    caso, parte da variabilidade total da resposta é "explicada" pelo bloco antes mesmo de
    comparar as técnicas, reduzindo o quadrado médio do erro. Se a gravidade basal \emph{não}
    estiver correlacionada com a resposta, blocar não introduz viés, mas também não reduz a
    variância — apenas consome $b-1=2$ graus de liberdade do erro sem retorno, exatamente como no
    exemplo do viveiro de pepineiros ($ER\approx1{,}02$) visto em aula.
\end{enumerate}
\end{solucao}

\begin{questao}{4} \textbf{(Agricultura --- dado faltante, fórmula de Yates)}
\end{questao}
\begin{solucao}
\begin{enumerate}[label=(\alph*)]
  \item Com $t=5$, $b=4$, $T'=58$, $B'=78$, $G'=380$:
    $$
    \hat y = \frac{t\,T' + b\,B' - G'}{(t-1)(b-1)} = \frac{5(58)+4(78)-380}{(4)(3)}
    = \frac{290+312-380}{12} = \frac{222}{12} = \textbf{18{,}5}.
    $$
  \item O DBCA completo (sem dado faltante) teria $gl_{Erro}=(t-1)(b-1)=(4)(3)=12$. Como o valor
    foi estimado (não observado), perde-se 1 grau de liberdade adicional:
    $gl_{Erro}^{ajustado} = 12-1=\textbf{11}$.
\end{enumerate}
\end{solucao}

\begin{questao}{5} \textbf{(Eficiência relativa)}
\end{questao}
\begin{solucao}
\begin{enumerate}[label=(\alph*)]
  \item $gl_{Bloco}=b-1=4$. $gl_{Erro}(DBCA)=(t-1)(b-1)=(5)(4)=20$.
    $gl_{Erro}(DCA)=gl_{Erro}(DBCA)+gl_{Bloco}=20+4=\textbf{24}$.
  \item
    $$
    ER = \frac{(gl_B+1)(gl_{E,DCA}+3)\,QM_{Bloco} + gl_{E,DCA}(gl_B+3)\,QM_{Erro}}
              {(gl_{E,DCA}+1)(gl_B+3)\,QM_{Erro}}
       = \frac{(5)(27)(45) + (24)(7)(15)}{(25)(7)(15)}
       = \frac{6075+2520}{2625} \approx \textbf{3{,}27}.
    $$
  \item Um DCA precisaria de aproximadamente $3{,}27$ vezes mais repetições por tratamento para
    atingir a mesma precisão deste DBCA — um ganho substancial, bem maior do que o
    $ER\approx1{,}02$ observado no experimento do viveiro de pepineiros. A diferença reflete o
    fato de que, aqui, $QM_{Bloco}=45$ é três vezes maior que $QM_{Erro}=15$ (os blocos capturam
    heterogeneidade real e considerável), enquanto no exemplo do pepino $QM_{Bloco}$ era menor que
    $QM_{Erro}$ — os blocos mal capturaram heterogeneidade alguma.
\end{enumerate}
\end{solucao}

\begin{questao}{6} \textbf{(Dissertativa --- matriz de delineamento)}
\end{questao}
\begin{solucao}
\begin{enumerate}[label=(\alph*)]
  \item A ausência de colunas de produto entre $X_{trat}$ e $X_{bloco}$ \textbf{é}, em termos do
    modelo, a suposição de \textbf{aditividade tratamento-bloco}: o efeito esperado de um
    tratamento não depende do bloco em que ele é observado (e vice-versa). Se houvesse uma
    interação real entre tratamento e bloco, o modelo $Y=X\beta+\varepsilon$ com essa matriz $X$
    estaria mal especificado, e a diferença ficaria absorvida (incorretamente) no termo de erro.
  \item Como as colunas de $X_{bloco}$ apenas identificam a qual bloco cada unidade pertence, e
    não carregam nenhuma informação de interesse por si — o "efeito de bloco" existe no modelo
    apenas para \textbf{absorver variação} que, de outra forma, cairia no erro —, testar se essa
    variação é "estatisticamente significativa" responde a uma pergunta diferente da que motivou
    incluir $X_{bloco}$ em $X$. Os blocos foram escolhidos pelo delineamento (não sorteados) com
    base em conhecimento prévio de heterogeneidade; removê-los do modelo por causa de um p-valor
    não significativo devolveria essa variação conhecida ao erro, sem nenhum ganho, e poderia até
    invalidar a estrutura de aleatorização restrita sob a qual os dados foram coletados.
\end{enumerate}
\end{solucao}

\begin{questao}{7} \textbf{(Ciência de dados --- fatores confundidos vs. blocagem)}
\end{questao}
\begin{solucao}
\begin{enumerate}[label=(\alph*)]
  \item O fator confundido é a \textbf{data de cadastro do usuário} (ou, mais amplamente,
    qualquer característica correlacionada a ela — tempo de uso do site, familiaridade com a
    interface, sazonalidade de quando o usuário chegou). Todo usuário cadastrado após 1º de
    janeiro de 2026 recebe \emph{sempre} o algoritmo novo, e todo usuário cadastrado antes recebe
    \emph{sempre} o antigo — as duas variáveis variam em perfeita sincronia.
  \item Como \texttt{Algoritmo} é uma função determinística da data de cadastro (um implica o
    outro sem exceção), as colunas indicadoras dessas duas variáveis em $X$ são idênticas a menos
    de uma constante — a matriz de delineamento não tem posto completo nessas duas colunas.
    \texttt{lm()} reportaria coeficiente \texttt{NA} para uma delas (tipicamente a última incluída
    na fórmula): não é possível separar estatisticamente o efeito do algoritmo do efeito de ser um
    usuário mais novo, \textbf{não importa quantos usuários adicionais sejam observados} sob essa
    mesma regra de atribuição — o problema é de delineamento, não de tamanho de amostra.
  \item Uma reformulação correta blocaria por \textbf{coorte de cadastro} (por exemplo, faixas de
    data de cadastro: usuários de 2023, 2024, 2025, 2026) — cada coorte funcionando como um bloco
    que agrupa usuários com veterania semelhante. \textbf{Dentro de cada bloco/coorte}, os
    usuários seriam então \textbf{sorteados} aleatoriamente entre algoritmo novo e antigo (um
    teste A/B tradicional restrito a cada coorte), de modo que ambos os algoritmos sejam
    observados em todo nível de veterania — quebrando a sincronia perfeita que causava a confusão
    no desenho original.
\end{enumerate}
\end{solucao}

\end{document}
